The low-noise amplifier sets the noise figure of the whole receiver (Friis), so
its design optimizes noise before gain.

\section{Noise Parameters}
A two-port's noise is fully described by four numbers: the minimum noise figure
$F_{min}$, the optimum source reflection $\Gamma_{opt}$ that achieves it, and the
equivalent noise resistance $R_n$ that says how quickly noise degrades away from
$\Gamma_{opt}$:
\[ F = F_{min} + \frac{4 R_n}{Z_0}\,\frac{|\Gamma_s - \Gamma_{opt}|^2}
       {(1-|\Gamma_s|^2)\,|1+\Gamma_{opt}|^2}. \]

\section{The Noise--Gain Trade}
In general the source impedance for minimum noise ($\Gamma_{opt}$) is \emph{not}
the one for maximum gain or best input match. Noise circles and gain circles on
the Smith chart make the compromise visible; a good LNA source match lands near
$\Gamma_{opt}$ and accepts slightly less than maximum gain.

\section{Inductive Source Degeneration}
The elegant trick used in CMOS/SiGe LNAs is a small source inductor $L_s$. It
generates a real input resistance $\approx \omega_T L_s$ \emph{without adding
thermal noise} of its own, so the input can be matched to $50\,\Omega$ at the same
time as the noise match --- the celebrated simultaneous noise-and-input match.

\section{Output and Bias}
With the source chosen for low noise, the output is conjugately matched for gain,
and a clean bias network keeps the stage stable. The result is a stage with a
noise figure within a few tenths of a dB of $F_{min}$ and enough gain to suppress
the noise of everything after it.
